% !TEX root = ../main.tex
\subsection{Dataset}
\label{sec:dataset}

\textsc{SciFig-Eval} comprises 250 English-language scientific figures from 187 arXiv papers (2023--2025), each paired with an expert description produced by two trained annotators (94\% agreement) with a third adjudicator for disagreements. \am{We selected bar charts, line plots, and pie charts because they were the most prevalent visualization types in the collected corpus, enabling reliable human annotation and evaluation. Although the benchmark begins with 250 figures, they are expanded through image transformations, reasoning questions, resistance probes, caption-bias probes, and selective-blur probes into over 34,000 evaluation instances across eight models.} Table~\ref{tab:dataset-summary} summarises the benchmark components; full details appear in Appendix~\ref{sec:dataset-details} (Table~\ref{tab:dataset-comprehensive}).

\begin{table}[t]
\centering
% \scriptsize
\footnotesize
\setlength{\tabcolsep}{3pt}
\begin{tabular}
% {@{}lr@{}}
{p{0.7\columnwidth}p{0.1\columnwidth}}
\toprule
\textbf{Component} & \textbf{Count} \\
\midrule
Baseline figures with expert descriptions (99B\,/\,99L\,/\,52P) & 250 \\
Figure transforms & 1{,}686 \\
Reasoning questions with human-reviewed solutions & 1{,}000 \\
Behavioural probes & 1{,}293 \\
\bottomrule
\end{tabular}
\caption{Benchmark summary. B\,=\,bar, L\,=\,line, P\,=\,pie.}
\label{tab:dataset-summary}
\end{table}
